\section{The row-column graph and general bounds}\label{sec:graph}

A blackout $S$ is \emph{axis-valid} when it distinguishes axis-aligned
rectangles. Call a rectangle \emph{weak} if at most two corners show.
In the \emph{row-column graph} $G_S$, join row $r_j$ to column $c_i$
for each black point $(i,j)$.

\begin{lemma}[axis-aligned rectangles and short cycles]\label{lem:girth}
If $G_S$ has no four-cycle and no six-cycle, then $S$ is axis-valid. The converse holds for $m\ge3$.
\end{lemma}

\begin{proof}
For $m\ge3$, a four-cycle $r_j\,c_i\,r_{j'}\,c_{i'}$ gives
indistinguishable rectangles $\{i,w\}\times\{j,j'\}$ and
$\{i',w\}\times\{j,j'\}$ using a third column $w$.
A six-cycle $r_a\,c_x\,r_b\,c_y\,r_c\,c_z$ gives
$\{x,y\}\times\{a,b\}$ and $\{y,z\}\times\{a,c\}$, both black
except possibly at $(y,a)$.
Conversely, axis-aligned $R_1\ne R_2$ with equal presentations have
$C(R_1)\triangle C(R_2)\subseteq S$. They cannot share opposite corners,
so by symmetry any common corners lie in one row $a$.
With none, $R_1$ is fully black. With two, $(x,a)$ and $(y,a)$,
the other rows $b\ne c$ give a black four-cycle $r_b\,c_x\,r_c\,c_y$.
With just $(x,a)$, write $R_1=\{x,y\}\times\{a,b\}$ and
$R_2=\{x,z\}\times\{a,c\}$ to obtain
$r_a\,c_y\,r_b\,c_x\,r_c\,c_z$.
\end{proof}

\begin{theorem}[row bound]\label{thm:rowbound}
Let $m\ge3$. Every axis-valid blackout of the $h\times m$ grid has at most $m+q_h$ points, and at most $h+q_m$ if also $h\ge3$. Hence $f(h,m)\le m+q_h$.
\end{theorem}

Join each nonempty column's black rows $B_i$ by a tree $T_i$.
Their union $H$ reduces the problem to Mantel's theorem.

\begin{lemma}[row graph]\label{lem:H}
Let $m\ge3$ and let $S$ be axis-valid. For any choice of the trees, $H$ is simple and triangle-free, and $e(H)=|S|-t$, where $t$ is the number of nonempty columns.
\end{lemma}

\begin{proof}
A repeated edge of $H$ would put the same row pair in two columns,
giving a four-cycle in $G_S$. Lemma~\ref{lem:girth} rules this out, so
$e(H)=\sum(|B_i|-1)=|S|-t$.
Now suppose $H$ contains a triangle $r_ar_br_c$. Its edges cannot all
belong to one tree. If two columns $i,i'$ supply them, two of $a,b,c$
belong to $B_i\cap B_{i'}$, again giving a four-cycle. If three columns
supply them, they give the six-cycle
$r_a\,c_i\,r_b\,c_{i'}\,r_c\,c_{i''}$. Both possibilities contradict
Lemma~\ref{lem:girth}.
\end{proof}

\begin{lemma}[Mantel \cite{mantel1907}]\label{lem:mantel}
Let $H$ be a simple triangle-free graph on $h$ vertices, let a vertex of maximum degree $\Delta$ have neighbor set $A$, and let $B$ be the other $h-\Delta$ vertices, spanning $e(B)$ edges. Then
\[
e(H)=\sum_{b\in B}d(b)-e(B)\le(h-\Delta)\Delta-e(B)\le q_h.
\]
\end{lemma}

\begin{proof}
The set $A$ is independent, so every edge meets $B$.
Summing degrees over $B$ counts its internal edges twice.
Subtract those once, then use $d(b)\le\Delta$ and
$(h-\Delta)\Delta\le q_h$.
\end{proof}

\begin{proof}[Proof of Theorem~\ref{thm:rowbound}]
Combining Lemmas~\ref{lem:H} and~\ref{lem:mantel} gives
$|S|=t+e(H)\le m+q_h$. The need for a third column enters only through
Lemma~\ref{lem:girth}. Transposing therefore gives $|S|\le h+q_m$ when
$h\ge3$.
\end{proof}

\begin{theorem}[equality structure]\label{thm:equality}
Let $h\ge4$ and $m\ge3$. A blackout of the $h\times m$ grid with $m+q_h$ points is axis-valid if and only if
\begin{enumerate}[(a)]
\item every column is a singleton or a double,
\item the $q_h$ doubles have distinct row pairs, and
\item these pairs form a balanced complete bipartite graph on the rows.
\end{enumerate}
There are $p_h\,(m)_{q_h}\,h^{m-q_h}$ such blackouts, where $(m)_q$ is the falling factorial and the number $p_h$ of balanced bipartitions of the rows is $\binom h{\lfloor h/2\rfloor}$, halved for even $h$.
\end{theorem}

\begin{proof}
Equality in the row-bound proof forces $t=m$, $e(B)=0$, and
$(h-\Delta)\Delta=q_h$, with all of $B$ joined to all of $A$.
Thus $H$ is the graph in (c).
If $|B_i|\ge3$, its tree $T_i$ has a vertex $v$ with neighbors $u,w$.
Since $h\ge4$, the part containing $v$ has another vertex $z$.
Replacing $vw$ by $uw$ preserves $uz,wz$ and creates triangle $uwz$,
contradicting Lemma~\ref{lem:H}. This proves (a), and the tree edges,
now the row pairs of doubles, give (b) and (c).
Conversely, (a) and (b) exclude four-cycles. A six-cycle
$r_a\,c_x\,r_b\,c_y\,r_c\,c_z$ needs doubles $\{a,b\}$,
$\{b,c\}$, and $\{a,c\}$, contradicting (c).
Apply Lemma~\ref{lem:girth}. To count, choose the bipartition,
assign edges to distinct columns, and choose the remaining singleton rows.
\end{proof}

\begin{corollary}[small widths]\label{cor:small-m}
Let $h\ge4$ and $3\le m<q_h$. Then $f(h,m)\le m+q_h-1$. In particular $f(5,5)\le10$.
\end{corollary}

\begin{proof}
Attaining $m+q_h$ would require $q_h\le m$ double columns by
Theorem~\ref{thm:equality}. Since this is impossible here,
Theorem~\ref{thm:rowbound} improves by one. Here
$q_5=6$.
\end{proof}

Put $W_h=1+\lfloor(h-1)^2/4\rfloor$, so $W_3=2$, $W_4=3$, and $W_5=5$.

\begin{lemma}[leftmost corner]\label{lem:leftmost}
A tilted rectangle $R$ has a unique corner $\ell=(x,r)$ in its smallest column, and a unique corner in its smallest row. The sides from $\ell$ are $(a,-b)$ and $(c,d)$ with positive integers $a,b,c,d$ and $ac=bd$, and the corners are
\begin{equation}\label{eq:tilt-param}
(x,r),\qquad(x+a,\,r-b),\qquad(x+c,\,r+d),\qquad(x+a+c,\,r-b+d).
\end{equation}
Its rows are $r-b$ to $r+d$, and its \emph{span} is $a+c$.
\end{lemma}

\begin{proof}
There are no vertical sides. Two corners in the leftmost column
$x$ would be opposite, forcing the other two there as well by their
mean column. Likewise for rows.
The two sides from $\ell$ therefore point right, one up and one down.
Perpendicularity gives $ac=bd$.
\end{proof}

\begin{lemma}[span bound]\label{lem:span}
Every tilted rectangle of the $h\times m$ grid has span $a+c\le W_h$.
\end{lemma}

\begin{proof}
The rectangle fits in $h$ rows, so $b+d\le h-1$. Also
$(a-1)(c-1)\ge0$, whence $a+c\le bd+1\le W_h$, using the integer bound
$bd\le(b+d)^2/4$.
\end{proof}

Thus $W_h$ is \seqnum{A033638}$(h-1)$. For $h\ge3$ and $m>W_h$,
the sides $(1,-b)$ and $(bd,d)$ attain it when $b+d=h-1$ and $|b-d|\le1$.
Here a rectangle's \emph{height} and \emph{width} mean its row and column
ranges (the distances between its extreme coordinates).

% TODO(author): check that this matches the wording of the definition and the conjecture in A338885.
\begin{corollary}[the conjecture in \seqnum{A338885}]\label{cor:oeisconj}
For $n\ge2$, the largest width of a tilted lattice rectangle of height $n$ is $\lfloor n^2/4\rfloor+1$. So row $n$ of \seqnum{A338885} has largest entry \seqnum{A033638}$(n)$.
\end{corollary}

\begin{proof}
In~\eqref{eq:tilt-param}, set the height $b+d=n$. As in the span bound,
the width is $a+c\le bd+1\le\lfloor n^2/4\rfloor+1$.
Choose sides $(1,-b)$ and $(bd,d)$ with $b=\lfloor n/2\rfloor$ and
$d=\lceil n/2\rceil$ to attain it.
\end{proof}

The \emph{diamond} $D(i,j)$ is the square with corners $(i,j+1)$, $(i+1,j)$, $(i+2,j+1)$, and $(i+1,j+2)$. The rectangles $\kA_i$, $\kB_i$, $\kC_i$, and $\kE_i$ have one corner in each of the columns $i$ to $i+3$, in the rows $3412$, $3142$, $2413$, and $2143$ from left to right.

\begin{lemma}[tilted rectangles on narrow grids]\label{lem:tilted-narrow}
The tilted rectangles of the $3\times m$ grid are the diamonds $D(i,1)$. Those of the $4\times m$ grid are $D(i,1)$, $D(i,2)$, $\kA_i$, $\kB_i$, $\kC_i$, and $\kE_i$.
\end{lemma}

\begin{proof}
All the listed sets are rectangles. On three rows, necessarily
$b=d=1$, hence $a=c=1$ and the leftmost corner $\ell$ is in row $2$.
On four rows, $b=d=1$ again gives diamonds. The only other possibility
is $\{b,d\}=\{1,2\}=\{a,c\}$, with $\ell$ in row $1+b$.
The four choices of sides, $(2,-2),(1,1)$, $(1,-2),(2,1)$,
$(2,-1),(1,2)$, and $(1,-1),(2,2)$, give $\kA$, $\kB$, $\kC$, and $\kE$.
\end{proof}

\begin{lemma}[boundary cross]\label{lem:cross}
The first row and the first column of the $h\times m$ grid form a valid blackout, so $f(h,m)\ge h+m-1$.
\end{lemma}

\begin{proof}
Only rectangles showing at most two points need checking
(Lemma~\ref{lem:threecorners}). An axis-aligned rectangle
$\{x_1,x_2\}\times\{y_1,y_2\}$ shows $X\times Y$, where
$X=\{x_1,x_2\}\setminus\{1\}$ and $Y=\{y_1,y_2\}\setminus\{1\}$.
These recover its columns and rows, since $|X|=1$ precisely when
$1\in\{x_1,x_2\}$, and likewise for $Y$.

By Lemma~\ref{lem:leftmost}, a tilted rectangle can have only two black
corners, $\ell$ and $(x+a,r-b)$. When both are black,
$\ell=(1,1+b)$ and the shown corners are $p=(1+c,1+b+d)$ and
$q=(1+a+c,1+d)$. Their coordinates recover $a,b,c,d$. They also differ
in both coordinates, so an axis-aligned rectangle could show this pair
only if its other two corners were black, which they are not.
\end{proof}

\begin{lemma}[axis-aligned rivals]\label{lem:rival}
Let no axis-aligned rectangle be fully black. A \emph{rival} of a set $P$ of at most two white points is an axis-aligned rectangle that shows exactly $P$. The empty set has no rival, and $P$ has one exactly when it is a singleton $\{(x,r)\}$ with $\{r,s\}\in F$ for some $s\in B_x$, a vertical pair $\{(x,r),(x,s)\}$ with $\{r,s\}\in F$, a horizontal pair $\{(x,r),(z,r)\}$ with $B_x\cap B_z\ne\emptyset$, or an oblique pair $\{(x,r),(z,s)\}$ with $x\ne z$, $r\ne s$, and $(x,s),(z,r)\in S$.
\end{lemma}

\begin{proof}
A rival of $P\ne\emptyset$ has rows $\{r,s\}$ and columns $\{x,z\}$,
with every corner outside $P$ black. Reading those corners gives the
listed conditions, which force $z\ne x$ and $s\ne r$ since $(x,r)$ is white.
\end{proof}
